\startSection
	\selectlanguage{french}
	\sectionTitle{L'Hymne National Français}
	
	\begin{question}{La Marseillaise, hymne emblématique de la République française, a été composée en 1792 dans un contexte de guerre révolutionnaire. Quel musicien est à l'origine de cette œuvre patriotique devenue symbole national ?}
		\choice[correct]{Claude Joseph Rouget de Lisle}
		\choice{Jean-Baptiste Lully}
		\choice{Hector Berlioz}
		\choice{François-Joseph Gossec}
		\choice{Jean-Philippe Rameau}
	\end{question}

	
	\begin{question}{La ville où fut composée la Marseillaise joue un rôle important dans son histoire et son nom. Dans quelle ville française ce chant révolutionnaire a-t-il été créé ?}
		\choice{Paris}
		\choice{Marseille}
		\choice{Lyon}
		\choice[correct]{Strasbourg}
		\choice{Bordeaux}
	\end{question}

	
	\begin{question}{Avant d’être connue sous le nom de \enquote{La Marseillaise}, cette œuvre portait un autre titre reflétant son usage militaire initial. Quel était le titre original donné à ce chant lors de sa création ?}
		\choice{Chant de guerre de l'armée française}
		\choice{Hymne des Marseillais}
		\choice[correct]{Chant de guerre pour l'armée du Rhin}
		\choice{Chant patriotique de la République}
		\choice{Marche des volontaires}
	\end{question}

	
	\begin{question}{La reconnaissance officielle d’un hymne national est un acte symbolique fort. En quelle année la Marseillaise a-t-elle été adoptée comme hymne national officiel de la France ?}
		\choice{1792}
		\choice[correct]{1795}
		\choice{1804}
		\choice{1815}
		\choice{1848}
	\end{question}

	
	\begin{question}{Le nom \enquote{La Marseillaise} ne provient pas de son lieu de composition, mais d’un événement marquant lié à sa diffusion. Pourquoi ce chant porte-t-il ce nom ?}
		\choice{Il a été composé à Marseille}
		\choice{Il célèbre une victoire à Marseille}
		\choice[correct]{Les fédérés marseillais l'ont popularisé à Paris}
		\choice{Le compositeur était marseillais}
		\choice{Il a été chanté pour la première fois à Marseille}
	\end{question}
\renderSection
